\documentclass{zhvt-classic}

% 使用本地源流明体（GenRyuMin2 TW）按文件路径加载，不依赖系统字体名
% 版框（内框／界栏／书口侧栏／鱼尾）改用与版心印章同一印泥棕色：
%   印章图片 yshs.png 实色 = RGB(135,68,52) = HTML 874434
\usepackage{xcolor}
\usepackage{tikz}
% sealbrown 双注册：xcolor 供 TikZ/pgf，L3 供 cls 版框
\definecolor{sealbrown}{HTML}{874434}
\ExplSyntaxOn
\color_set:nnn { sealbrown } { HTML } { 874434 }
\ExplSyntaxOff

\zhvtset{
  grid color       = sealbrown,
  main font        = GenRyuMin2TW-SB,
  main font path   = C:/Users/Administrator/texmf/fonts/opentype/public/GenRyuMin2TW/,
jiazhu font      = GenRyuMin2TW-M,
  jiazhu font path = C:/Users/Administrator/texmf/fonts/opentype/public/GenRyuMin2TW/,
  % 页码框容纳 4 字：正常页「九十七」+前置空 1 格=4 行恰好；
  % 100+ 页「一百零一」共 5 行超出，末字「一/二/三」被框裁掉（2026-10-04 实测），故放宽到 5
  shukou page number chars = 5,
  % 书口书名框高从默认 6 字减到 4 字，把省出的 2 字高拨给 chapter 框，
  % 使 6 字书名+卷名放进 9 字高的 chapter 框后不再压线（2026-10-05）
  shukou title chars = 4,
}

\title{史记}[卫将军骠骑列传]
\author{汉·司马迁}
\maker{闲叟漫记·重製}


% ===== 索隐述赞样式（2026-11-08 四次定稿 · 愚大人亲手改定参数）=====
% ⚠ 竖排前提：自定义字体家族必须带 RawFeature={vertical:+vert}；
%   simhei.ttf 加 +vhal 会整体加载失败静默回退，只写 +vert；simfang 同理。
\newCJKfontfamily{\shuzanfont}[Path=C:/Windows/Fonts/,RawFeature={vertical:+vert}]{simhei.ttf}
\newCJKfontfamily{\shuzanbody}[Path=C:/Windows/Fonts/,RawFeature={vertical:+vert}]{simfang.ttf}
% 题头：152x30pt 圆角 6pt 棕底白字；baseline (0,2) 使整框落入界行格子不压版框线；
% 白字节点 x 不补偿、y +8.2pt（像素级墨迹实测 dx/dy≈0，L/R 留白约 20pt）。
\newcommand{\shuzanhead}[1]{%
  \begin{tikzpicture}[baseline={(0,2)}]
    \fill[sealbrown, rounded corners=6pt]
      (-76pt,-15pt) rectangle (76pt,15pt);
    \node[text=white, inner sep=0pt, anchor=center] at (0pt,8.2pt)
      {\shuzanfont\fontsize{23}{23}\selectfont #1};
  \end{tikzpicture}%
}
% ★ 修正书口卷名位置：chapter 框起始位置改回 cls 默认值，避免全局影响。
%   只在每个卷首页通过 \juan 命令里的空命令微调该页位置（安全，不会跨页污染）。
\ExplSyntaxOn
\AtBeginDocument{
  \dim_set:Nn \l__zhvt_shukou_chapter_title_start_dim
    { \l__zhvt_shukou_title_height_dim + \l__zhvt_font_size_dim * 2 }
  \dim_set:Nn \l__zhvt_shukou_page_number_start_dim
    { \l__zhvt_shukou_chapter_title_start_dim + \l__zhvt_shukou_chapter_title_height_dim }
}
\ExplSyntaxOff

\begin{document}

% ★★ 版心（书口）排布覆写：把书名放到鱼尾之下（2026-10-04）
%
%   cls 默认版心自上而下：书名（title 静态框，在鱼尾之上）→ 鱼尾 → 篇名（chapter 动态框）→ 页码
%   本模板改为：          鱼尾（位置不动）→ 书名 → 篇名 → 页码分割线 → 页码
%
%   手法：把 cls 放在 title 静态框里的书名，改放进 chapter 动态框（鱼尾之下）；
%         title 静态框留空。书名取 \l__zhvt_title_str，即 \title{...} 的第一参数，
%         所以换书时只改 \title，不必在此重复写书名。
%
%   ⚠ 三条必须遵守的写法（2026-10-04 实测）：
%     ① 汉字必须用 \hbox{...} 包住。flowfram 会把框内容装进一个宽 14.4mm（= 一个 line_sep
%        = 书口栏宽）的纵向 minipage 并按该宽度折行，而整页靠 PDF /Rotate 90 显示、
%        框内文字本身不变向 —— 不包 \hbox 时「围炉夜话」4 字会被逐字折成 4 行，
%        在版面上就成了横排 4 字并横向溢出版框。
%     ② 竖排下**一个 \hbox 就是一整列**。书名与篇名要同处一列，必须放进同一个 \hbox；
%        分成两个 \hbox 时，后一个会被当成新的一列、甩到隔壁栏去
%        （实测「序」跑到书口栏外侧 14.4mm 的正文列上）。中间的 \quad 在框内是横向空白，
%        显示出来正好是书名与篇名之间的纵向间隔。
%     ③ ★ 内容盒要写成 **与框等宽的定宽盒** `\hbox to \l__zhvt_shukou_chapter_title_height_dim{…\hfil}`。
%        实测：chapter 框是把内容盒沿纵向**居中**摆放的，内容宽度一变，书名就跟着漂：
%        自然宽度的 \hbox（书名 38.4mm）落在 x=157.2mm，加篇名后（57.6mm）落在 x=152.4mm，
%        目录页（无篇名）与正文页的书名就差半格。做成等宽盒后，无论有没有篇名，
%        内容固定从框顶起排，实测两页均落在 x=142.8mm，不再漂。
\ExplSyntaxOn
% 把书口 chapter 框（书名+篇名）起始位置下移 1 字高，避免 4 字卷名顶到鱼尾
% ⚠ cls 的 \__zhvt_calculate_dimensions: 在 \AtBeginDocument 里执行，会重算该 dim；
%   必须把覆写挂到 begin document 之后，否则会被二次覆盖（2026-10-05 实测）。
% (旧方案已移到导言区，直接修改 chapter 框起始位置)
% 卷标题命令 #1=书口篇名(鱼尾) #2=正文卷名（含卷序，卷名内空格用 ~）
%
% ★★ 目录页两项修订（2026-10-05，愚大人指出）★★
%   ① 页码起点统一对齐：
%      原状 —— 目录项写成 \par\hspace{4em}卷名\hspace{4em}页码，页码起点 = 卷名实宽 + 4em。
%      卷名实宽随字数变：「卷一~梁惠王上」6 字 +~、「卷七~离娄上」5 字 +~，
%      差 1 个字位 → 卷七~卷十 4 条的页码整体比其余 10 条高 27.21pt（一个字），看着上下不齐。
%      改法 —— 卷名补全角空格（宽 = 1 字 = 27.215pt）补足到 **7 字符位**再写进目录，
%      14 条页码起点遂全部落在同一条竖线上（起点对齐；末字随页码字数自然参差）。
%      ⚠ 只对目录里的副本 \l_tmpc_tl 补，正文卷名仍用 \l_tmpa_tl 原样 —— 正文像素不受影响。
%   ② 卷名整体上移半字位：前导 \hspace 由 4em 改 3.5em。
%      实测原状目录块在版框内偏下（天头留白 139.6pt / 地脚留白 110.7pt），
%      上移 0.5em = 13.6pt 后上下基本均分，视觉居中。
\NewDocumentCommand { \juan } { mm } {
	\clearpage
	\tl_set:Nn \l__zhvt_chapter_title_tl { #1 }
	\tl_set:Nn \l_tmpa_tl { #2 }
	\tl_set_eq:NN \l_tmpc_tl \l_tmpa_tl
	\int_step_inline:nn { 7 - \str_count:n { #2 } }
	  { \tl_put_right:Nn \l_tmpc_tl { 　 } }
	\tl_set:Nx \l_tmpb_tl { \zhnum { page } }
	\tl_put_left:Nn \l_tmpb_tl { \hspace { 4em } }
	\tl_put_left:Nn \l_tmpb_tl { \l_tmpc_tl }
	\tl_put_left:Nn \l_tmpb_tl { \par \hspace { 3.5em } }
	\addtocontents { toc } { \l_tmpb_tl }
	\skip_horizontal:n { \l__zhvt_font_size_dim * 4 }
	\l_tmpa_tl
	\par
	\thispagestyle { empty }
}
\cs_set_protected:Npn \__zhvt_set_shukou_content: {
	% ★ 鱼尾上方留白（title 静态框）嵌入闲叟钤印：
	%   整页 /Rotate 90，故此处把图用 \rotatebox{90} 转正为竖章；
	%   尺寸卡在版心条幅内（书内约 7mm 宽 × 33.8mm 高，留余量不压鱼尾）。
	\setstaticcontents * { title } {
		\vbox_to_ht:nn { \l__zhvt_baseline_skip_dim } {
			\vskip 14.2pt
			\hbox_to_wd:nn { \l__zhvt_shukou_title_height_dim } {
				\hfil
				\rotatebox [ origin = c ] { 90 } { \includegraphics [ width = 7mm , height = 33.8mm ] { yshs.png } }
				\hfil
			}
		}
	}
% ★ 版心（书口）统一改法（2026-11-16 愚大人定）：18pt＋框内居中＋书口文字同版框色
%   ① 原 27.2pt 时 7 字长篇名（如「廉颇蔺相如列传」）末字溢出 chapter 框、压住页码；
%      改 \fontsize{18pt}{18pt}，前后 \hfil 令字串在框内居中（不随篇名字数漂移）。
%   ② \vbox{\vskip 8.2pt} 定量下移，抵消缩字号后栏宽方向 8.2pt 上偏（像素实测 y 中心 543.3 = 原版）。
%      ⚠ \raisebox / \vspace* 在 flowfram 框内无效；换字号须同步换补偿：16pt→10.1、19pt→7.4。
%   ③ chapter 与 page~number 均置 \color{sealbrown}（HTML 874434），与版框／界行同色；正文仍黑。
%   ⚠ 页码字位、钤印、鱼尾、界行、正文一律未动；五部合订本同此口径。
	\setdynamiccontents * { chapter }
	{ \color{sealbrown} \vbox { \vskip 8.2pt \hbox to \l__zhvt_shukou_chapter_title_height_dim
		{ \hfil \fontsize{18pt}{18pt}\selectfont \l__zhvt_title_str \quad \l__zhvt_chapter_title_tl \hfil } } }
	\setdynamiccontents * { page~number } { \color{sealbrown} \quad \zhnum { page } }
}
\__zhvt_set_shukou_content:
\ExplSyntaxOff

\maketitle

% 单篇，无多卷目录，不渲染目录页

\juan{卫将军骠骑列传}{卷一百十一~卫将军骠骑列传第五十一}
\hspace*{2em}大将军卫青者。平阳人也。其父郑季。为吏。给事平阳侯家。与侯妾卫媪通。生青。青同母兄卫长子。而姊卫子夫自平阳公主家得幸天子。故冒姓为卫氏。字仲卿。长子更字长君。长君母号为卫媪。媪长女卫孺。次女少兒。次女卫子夫。后子夫男弟步。广皆冒卫氏。\\
\hspace*{2em}青为侯家人。少时归其父。其父使牧羊。先母之子皆奴畜之。不以为兄弟数。青尝从入至甘泉居室。有一钳徒相青曰。﹁贵人也。官至封侯。﹂青笑曰。﹁人奴之生。得毋笞骂即足矣。安得封侯事乎。﹂\\
\hspace*{2em}青壮。为侯家骑。从平阳主。建元二年春。青姊子夫得入宫幸上。皇后。堂邑大长公主女也。无子。妒。大长公主闻卫子夫幸。有身。妒之。乃使人捕青。青时给事建章。未知名。大长公主执囚青。欲杀之。其友骑郎公孙敖与壮士往篡取之。以故得不死。上闻。乃召青为建章监。侍中。及同母昆弟贵。赏赐数日间累千金。孺为太仆公孙贺妻。少兒故与陈掌通。上召贵掌。公孙敖由此益贵。子夫为夫人。青为大中大夫。\\
\hspace*{2em}元光五年。青为车骑将军。击匈奴。出上谷。太仆公孙贺为轻车将军。出云中。大中大夫公孙敖为骑将军。出代郡。卫尉李广为骁骑将军。出雁门。军各万骑。青至茏城。斩首虏数百。骑将军敖亡七千骑。卫尉李广为虏所得。得脱归。皆当斩。赎为庶人。贺亦无功。\\
\hspace*{2em}元朔元年春。卫夫人有男。立为皇后。其秋。青为车骑将军。出雁门。三万骑击匈奴。斩首虏数千人。明年。匈奴入杀辽西太守。虏略渔阳二千馀人。败韩将军军。汉令将军李息击之。出代。令车骑将军青出云中以西至高阙。遂略河南地。至于陇西。捕首虏数千。畜数十万。走白羊。楼烦王。遂以河南地为朔方郡。以三千八百户封青为长平侯。青校尉苏建有功。以千一百户封建为平陵侯。使建筑朔方城。青校尉张次公有功。封为岸头侯。天子曰。﹁匈奴逆天理。乱人伦。暴长虐老。以盗窃为务。行诈诸蛮夷。造谋藉兵。数为边害。故兴师遣将。以征厥罪。诗不云乎。﹁薄伐玁狁。至于太原﹂。﹁出车彭彭。城彼朔方﹂。今车骑将军青度西河至高阙。获首虏二千三百级。车辎畜产毕收为卤。已封为列侯。遂西定河南地。按榆谿旧塞。绝梓领。梁北河。讨蒲泥。破符离。斩轻锐之卒。捕伏听者三千七十一级。执讯获丑。驱马牛羊百有馀万。全甲兵而还。益封青三千户。﹂其明年。匈奴入杀代郡太守友。入略雁门千馀人。其明年。匈奴大入代。定襄。上郡。杀略汉数千人。\\
\hspace*{2em}其明年。元朔之五年春。汉令车骑将军青将三万骑。出高阙。卫尉苏建为游击将军。左内史李沮为彊弩将军。太仆公孙贺为骑将军。代相李蔡为轻车将军。皆领属车骑将军。俱出朔方。大行李息。岸头侯张次公为将军。出右北平。咸击匈奴。匈奴右贤王当卫青等兵。以为汉兵不能至此。饮醉。汉兵夜至。围右贤王。右贤王惊。夜逃。独与其爱妾一人壮骑数百驰。溃围北去。汉轻骑校尉郭成等逐数百里。不及。得右贤裨王十馀人。众男女万五千馀人。畜数千百万。於是引兵而还。至塞。天子使使者持大将军印。即军中拜车骑将军青为大将军。诸将皆以兵属大将军。大将军立号而归。天子曰。﹁大将军青躬率戎士。师大捷。获匈奴王十有馀人。益封青六千户。﹂而封青子伉为宜春侯。青子不疑为阴安侯。青子登为发干侯。青固谢曰。﹁臣幸得待罪行间。赖陛下神灵。军大捷。皆诸校尉力战之功也。陛下幸已益封臣青。臣青子在襁褓中。未有勤劳。上幸列地封为三侯。非臣待罪行间所以劝士力战之意也。伉等三人何敢受封。﹂天子曰。﹁我非忘诸校尉功也。今固且图之。﹂乃诏御史曰。﹁护军都尉公孙敖三从大将军击匈奴。常护军。傅校获王。以千五百户封敖为合骑侯。都尉韩说从大将军出窳浑。至匈奴右贤王庭。为麾下搏战获王。以千三百户封说为龙嵒侯。骑将军公孙贺从大将军获王。以千三百户封贺为南窌侯。轻车将军李蔡再从大将军获王。以千六百户封蔡为乐安侯。校尉李朔。校尉赵不虞。校尉公孙戎奴。各三从大将军获王。以千三百户封朔为涉轵侯。以千三百户封不虞为随成侯。以千三百户封戎奴为从平侯。将军李沮。李息及校尉豆如意有功。赐爵关内侯。食邑各三百户。﹂其秋。匈奴入代。杀都尉硃英。\\
\hspace*{2em}其明年春。大将军青出定襄。合骑侯敖为中将军。太仆贺为左将军。翕侯赵信为前将军。卫尉苏建为右将军。郎中令李广为后将军。右内史李沮为彊弩将军。咸属大将军。斩首数千级而还。月馀。悉复出定襄击匈奴。斩首虏万馀人。右将军建。前将军信并军三千馀骑。独逢单于兵。与战一日馀。汉兵且尽。前将军故胡人。降为翕侯。见急。匈奴诱之。遂将其馀骑可八百。饹降单于。右将军苏建尽亡其军。独以身得亡去。自归大将军。大将军问其罪正闳。长史安。议郎周霸等。﹁建当云何。﹂霸曰。﹁自大将军出。未尝斩裨将。今建弃军。可斩以明将军之威。﹂闳。安曰。﹁不然。兵法﹁小敌之坚。大敌之禽也﹂。今建以数千当单于数万。力战一日馀。士尽。不敢有二心。自归。自归而斩之。是示后无反意也。不当斩。﹂大将军曰。﹁青幸得以肺腑待罪行间。不患无威。而霸说我以明威。甚失臣意。且使臣职虽当斩将。以臣之尊宠而不敢自擅专诛於境外。而具归天子。天子自裁之。於是以见为人臣不敢专权。不亦可乎。﹂军吏皆曰﹁善﹂。遂囚建诣行在所。入塞罢兵。\\
\hspace*{2em}是岁也。大将军姊子霍去病年十八。幸。为天子侍中。善骑射。再从大将军。受诏与壮士。为剽姚校尉。与轻勇骑八百直弃大军数百里赴利。斩捕首虏过当。於是天子曰。﹁剽姚校尉去病斩首虏二千二十八级。及相国。当户。斩单于大父行籍若侯产。生捕季父罗姑比。再冠军。以千六百户封去病为冠军侯。上谷太守郝贤四从大将军。捕斩首虏二千馀人。以千一百户封贤为众利侯。﹂是岁。失两将军军。亡翕侯。军功不多。故大将军不益封。右将军建至。天子不诛。赦其罪。赎为庶人。\\
\hspace*{2em}大将军既还。赐千金。是时王夫人方幸於上。甯乘说大将军曰。﹁将军所以功未甚多。身食万户。三子皆为侯者。徒以皇后故也。今王夫人幸而宗族未富贵。原将军奉所赐千金为王夫人亲寿。﹂大将军乃以五百金为寿。天子闻之。问大将军。大将军以实言。上乃拜甯乘为东海都尉。\\
\hspace*{2em}张骞从大将军。以尝使大夏。留匈奴中久。导军。知善水草处。军得以无饥渴。因前使绝国功。封骞博望侯。\\
\hspace*{2em}冠军侯去病既侯三岁。元狩二年春。以冠军侯去病为骠骑将军。将万骑出陇西。有功。天子曰。﹁骠骑将军率戎士逾乌盭。讨遫濮。涉狐奴。历五王国。辎重人众慑慴者弗取。冀获单于子。转战六日。过焉支山千有馀里。合短兵。杀折兰王。斩卢胡王。诛全甲。执浑邪王子及相国。都尉。首虏八千馀级。收休屠祭天金人。益封去病二千户。﹂\\
\hspace*{2em}其夏。骠骑将军与合骑侯敖俱出北地。异道。博望侯张骞。郎中令李广俱出右北平。异道。皆击匈奴。郎中令将四千骑先至。博望侯将万骑在后至。匈奴左贤王将数万骑围郎中令。郎中令与战二日。死者过半。所杀亦过当。博望侯至。匈奴兵引去。博望侯坐行留。当斩。赎为庶人。而骠骑将军出北地。已遂深入。与合骑侯失道。不相得。骠骑将军逾居延至祁连山。捕首虏甚多。天子曰。﹁骠骑将军逾居延。遂过小月氏。攻祁连山。得酋涂王。以众降者二千五百人。斩首虏三万二百级。获五王。五王母。单于阏氏。王子五十九人。相国。将军。当户。都尉六十三人。师大率减什三。益封去病五千户。赐校尉从至小月氏爵左庶长。鹰击司马破奴再从骠骑将军斩遬濮王。捕稽沮王。千骑将得王。王母各一人。王子以下四十一人。捕虏三千三百三十人。前行捕虏千四百人。以千五百户封破奴为从骠侯。校尉句王高不识。从骠骑将军捕呼于屠王王子以下十一人。捕虏千七百六十八人。以千一百户封不识为宜冠侯。校尉仆多有功。封为煇渠侯。﹂合骑侯敖坐行留不与骠骑会。当斩。赎为庶人。诸宿将所将士马兵亦不如骠骑。骠骑所将常选。然亦敢深入。常与壮骑先其大军。军亦有天幸。未尝困绝也。然而诸宿将常坐留落不遇。由此骠骑日以亲贵。比大将军。\\
\hspace*{2em}其秋。单于怒浑邪王居西方数为汉所破。亡数万人。以骠骑之兵也。单于怒。欲召诛浑邪王。浑邪王与休屠王等谋欲降汉。使人先要边。是时大行李息将城河上。得浑邪王使。即驰传以闻。天子闻之。於是恐其以诈降而袭边。乃令骠骑将军将兵往迎之。骠骑既渡河。与浑邪王众相望。浑邪王裨将见汉军而多欲不降者。颇遁去。骠骑乃驰入与浑邪王相见。斩其欲亡者八千人。遂独遣浑邪王乘传先诣行在所。尽将其众渡河。降者数万。号称十万。既至长安。天子所以赏赐者数十巨万。封浑邪王万户。为漯阴侯。封其裨王呼毒尼为下摩侯。鹰庇为煇渠侯。禽犁为河綦侯。大当户铜离为常乐侯。於是天子嘉骠骑之功曰。﹁骠骑将军去病率师攻匈奴西域王浑邪。王及厥众萌咸相饹。率以军粮接食。并将控弦万有馀人。诛獟駻。获首虏八千馀级。降异国之王三十二人。战士不离伤。十万之众咸怀集服。仍与之劳。爰及河塞。庶几无患。幸既永绥矣。以千七百户益封骠骑将军。﹂减陇西。北地。上郡戍卒之半。以宽天下之繇。\\
\hspace*{2em}居顷之。乃分徙降者边五郡故塞外。而皆在河南。因其故俗。为属国。其明年。匈奴入右北平。定襄。杀略汉千馀人。\\
\hspace*{2em}其明年。天子与诸将议曰。﹁翕侯赵信为单于画计。常以为汉兵不能度幕轻留。今大发士卒。其势必得所欲。﹂是岁元狩四年也。\\
\hspace*{2em}元狩四年春。上令大将军青。骠骑将军去病将各五万骑。步兵转者踵军数十万。而敢力战深入之士皆属骠骑。骠骑始为出定襄。当单于。捕虏言单于东。乃更令骠骑出代郡。令大将军出定襄。郎中令为前将军。太仆为左将军。主爵赵食其为右将军。平阳侯襄为后将军。皆属大将军。兵即度幕。人马凡五万骑。与骠骑等咸击匈奴单于。赵信为单于谋曰。﹁汉兵既度幕。人马罢。匈奴可坐收虏耳。﹂乃悉远北其辎重。皆以精兵待幕北。而適值大将军军出塞千馀里。见单于兵陈而待。於是大将军令武刚车自环为营。而纵五千骑往当匈奴。匈奴亦纵可万骑。会日且入。大风起。沙砾击面。两军不相见。汉益纵左右翼绕单于。单于视汉兵多。而士马尚彊。战而匈奴不利。薄莫。单于遂乘六。壮骑可数百。直冒汉围西北驰去。时已昏。汉匈奴相纷挐。杀伤大当。汉军左校捕虏言单于未昏而去。汉军因发轻骑夜追之。大将军军因随其后。匈奴兵亦散走。迟明。行二百馀里。不得单于。颇捕斩首虏万馀级。遂至窴颜山赵信城。得匈奴积粟食军。军留一日而还。悉烧其城馀粟以归。\\
\hspace*{2em}大将军之与单于会也。而前将军广。右将军食其军别从东道。或失道。后击单于。大将军引还过幕南。乃得前将军。右将军。大将军欲使使归报。令长史簿责前将军广。广自杀。右将军至。下吏。赎为庶人。大将军军入塞。凡斩捕首虏万九千级。\\
\hspace*{2em}是时匈奴众失单于十馀日。右谷蠡王闻之。自立为单于。单于后得其众。右王乃去单于之号。\\
\hspace*{2em}骠骑将军亦将五万骑。车重与大将军军等。而无裨将。悉以李敢等为大校。当裨将。出代。右北平千馀里。直左方兵。所斩捕功已多大将军。军既还。天子曰。﹁骠骑将军去病率师。躬将所获荤粥之士。约轻赍。绝大幕。涉获章渠。以诛比车耆。转击左大将。斩获旗鼓。历涉离侯。济弓闾。获屯头王。韩王等三人。将军。相国。当户。都尉八十三人。封狼居胥山。禅於姑衍。登临翰海。执卤获丑七万有四百四十三级。师率减什三。取食於敌。逴行殊远而粮不绝。以五千八百户益封骠骑将军。﹂右北平太守路博德属骠骑将军。会与城。不失期。从至檮余山。斩首捕虏二千七百级。以千六百户封博德为符离侯。北地都尉邢山从骠骑将军获王。以千二百户封山为义阳侯。故归义因淳王复陆支。楼专王伊即靬皆从骠骑将军有功。以千三百户封复陆支为壮侯。以千八百户封伊即靬为众利侯。从骠侯破奴。昌武侯安稽从骠骑有功。益封各三百户。校尉敢得旗鼓。为关内侯。食邑二百户。校尉自为爵大庶长。军吏卒为官。赏赐甚多。而大将军不得益封。军吏卒皆无封侯者。\\
\hspace*{2em}两军之出塞。塞阅官及私马凡十四万匹。而复入塞者不满三万匹。乃益置大司马位。大将军。骠骑将军皆为大司马。定令。令骠骑将军秩禄与大将军等。自是之后。大将军青日退。而骠骑日益贵。举大将军故人门下多去事骠骑。辄得官爵。唯任安不肯。\\
\hspace*{2em}骠骑将军为人少言不泄。有气敢任。天子尝欲教之孙吴兵法。对曰。﹁顾方略何如耳。不至学古兵法。﹂天子为治第。令骠骑视之。对曰。﹁匈奴未灭。无以家为也。﹂由此上益重爱之。然少而侍中。贵。不省士。其从军。天子为遣太官赍数十乘。既还。重车馀弃粱肉。而士有饥者。其在塞外。卒乏粮。或不能自振。而骠骑尚穿域蹋鞠。事多此类。大将军为人仁善退让。以和柔自媚於上。然天下未有称也。\\
\hspace*{2em}骠骑将军自四年军后三年。元狩六年而卒。天子悼之。发属国玄甲军。陈自长安至茂陵。为冢象祁连山。谥之。并武与广地曰景桓侯。子嬗代侯。嬗少。字子侯。上爱之。幸其壮而将之。居六岁。元封元年。嬗卒。谥哀侯。无子。绝。国除。\\
\hspace*{2em}自骠骑将军死后。大将军长子宜春侯伉坐法失侯。后五岁。伉弟二人。阴安侯不疑及发干侯登皆坐酎金失侯。失侯后二岁。冠军侯国除。其后四年。大将军青卒。谥为烈侯。子伉代为长平侯。\\
\hspace*{2em}自大将军围单于之后。十四年而卒。竟不复击匈奴者。以汉马少。而方南诛两越。东伐朝鲜。击羌。西南夷。以故久不伐胡。\\
\hspace*{2em}大将军以其得尚平阳长公主故。长平侯伉代侯。六岁。坐法失侯。\\
\hspace*{2em}左方两大将军及诸裨将名。\\
\hspace*{2em}最大将军青。凡七出击匈奴。斩捕首虏五万馀级。一与单于战。收河南地。遂置朔方郡。再益封。凡万一千八百户。封三子为侯。侯千三百户。并之。万五千七百户。其校尉裨将以从大将军侯者九人。其裨将及校尉已为将者十四人。为裨将者曰李广。自有传。无传者曰。\\
\hspace*{2em}将军公孙贺。贺。义渠人。其先胡种。贺父浑邪。景帝时为平曲侯。坐法失侯。贺。武帝为太子时舍人。武帝立八岁。以太仆为轻车将军。军马邑。后四岁。以轻车将军出云中。后五岁。以骑将军从大将军有功。封为南窌侯。后一岁。以左将军再从大将军出定襄。无功。后四岁。以坐酎金失侯。后八岁。以浮沮将军出五原二千馀里。无功。后八岁。以太仆为丞相。封葛绎侯。贺七为将军。出击匈奴无大功。而再侯。为丞相。坐子敬声与阳石公主奸。为巫蛊。族灭。无后。\\
\hspace*{2em}将军李息。郁郅人。事景帝。至武帝立八岁。为材官将军。军马邑。后六岁。为将军。出代。后三岁。为将军。从大将军出朔方。皆无功。凡三为将军。其后常为大行。\\
\hspace*{2em}将军公孙敖。义渠人。以郎事武帝。武帝立十二岁。为骑将军。出代。亡卒七千人。当斩。赎为庶人。后五岁。以校尉从大将军有功。封为合骑侯。后一岁。以中将军从大将军。再出定襄。无功。后二岁。以将军出北地。后骠骑期。当斩。赎为庶人。后二岁。以校尉从大将军。无功。后十四岁。以因杅将军筑受降城。七岁。复以因杅将军再出击匈奴。至余吾。亡士卒多。下吏。当斩。诈死。亡居民间五六岁。后发觉。复系。坐妻为巫蛊。族。凡四为将军。出击匈奴。一侯。\\
\hspace*{2em}将军李沮。云中人。事景帝。武帝立十七岁。以左内史为彊弩将军。后一岁。复为彊弩将军。\\
\hspace*{2em}将军李蔡。成纪人也。事孝文帝。景帝。武帝。以轻车将军从大将军有功。封为乐安侯。已为丞相。坐法死。\\
\hspace*{2em}将军张次公。河东人。以校尉从卫将军青有功。封为岸头侯。其后太后崩。为将军。军北军。后一岁。为将军。从大将军。再为将军。坐法失侯。次公父隆。轻车武射也。以善射。景帝幸近之也。\\
\hspace*{2em}将军苏建。杜陵人。以校尉从卫将军青。有功。为平陵侯。以将军筑朔方。后四岁。为游击将军。从大将军出朔方。后一岁。以右将军再从大将军出定襄。亡翕侯。失军。当斩。赎为庶人。其后为代郡太守。卒。冢在大犹乡。\\
\hspace*{2em}将军赵信。以匈奴相国降。为翕侯。武帝立十七岁。为前将军。与单于战。败。降匈奴。\\
\hspace*{2em}将军张骞。以使通大夏。还。为校尉。从大将军有功。封为博望侯。后三岁。为将军。出右北平。失期。当斩。赎为庶人。其后使通乌孙。为大行而卒。冢在汉中。\\
\hspace*{2em}将军赵食其。祋祤人也。武帝立二十二岁。以主爵为右将军。从大将军出定襄。迷失道。当斩。赎为庶人。\\
\hspace*{2em}将军曹襄。以平阳侯为后将军。从大将军出定襄。襄。曹参孙也。\\
\hspace*{2em}将军韩说。弓高侯庶孙也。以校尉从大将军有功。为龙嵒侯。坐酎金失侯。元鼎六年。以待诏为横海将军。击东越有功。为按道侯。以太初三年为游击将军。屯於五原外列城。为光禄勋。掘蛊太子宫。卫太子杀之。\\
\hspace*{2em}将军郭昌。云中人也。以校尉从大将军。元封四年。以太中大夫为拔胡将军。屯朔方。还击昆明。毋功。夺印。\\
\hspace*{2em}将军荀彘。太原广武人。以御见。侍中。为校尉。数从大将军。以元封三年为左将军击朝鲜。毋功。以捕楼船将军坐法死。\\
\hspace*{2em}最骠骑将军去病。凡六出击匈奴。其四出以将军。斩捕首虏十一万馀级。及浑邪王以众降数万。遂开河西酒泉之地。西方益少胡寇。四益封。凡万五千一百户。其校吏有功为侯者凡六人。而后为将军二人。\\
\hspace*{2em}将军路博德。平州人。以右北平太守从骠骑将军有功。为符离侯。骠骑死后。博德以卫尉为伏波将军。伐破南越。益封。其后坐法失侯。为彊弩都尉。屯居延。卒。\\
\hspace*{2em}将军赵破奴。故九原人。尝亡入匈奴。已而归汉。为骠骑将军司马。出北地时有功。封为从骠侯。坐酎金失侯。后一岁。为匈河将军。攻胡至匈河水。无功。后二岁。击虏楼兰王。复封为浞野侯。后六岁。为浚稽将军。将二万骑击匈奴左贤王。左贤王与战。兵八万骑围破奴。破奴生为虏所得。遂没其军。居匈奴中十岁。复与其太子安国亡入汉。后坐巫蛊。族。\\
\hspace*{2em}自卫氏兴。大将军青首封。其后枝属为五侯。凡二十四岁而五侯尽夺。卫氏无为侯者。\\
\hspace*{2em}太史公曰。苏建语余曰。﹁吾尝责大将军至尊重。而天下之贤大夫毋称焉。原将军观古名将所招选择贤者。勉之哉。大将军谢曰。﹁自魏其。武安之厚宾客。天子常切齿。彼亲附士大夫。招贤绌不肖者。人主之柄也。人臣奉法遵职而已。何与招士。﹂﹂骠骑亦放此意。其为将如此。\\

% ===== 索隐述赞（司马贞《史记索隐》）· 统一样式 =====
% 题头「索隐述赞曰」五字并入反白框；正文每四句一列，定宽盒 \hfil 居列中；
% \fontsize{24}{41} 列距对齐全书界行（121+41k）；\clearpage 使述赞独占一页。
\clearpage
\shuzanhead{索隐述赞曰}\linebreak
{\shuzanbody\fontsize{24}{41}\selectfont
\hbox to \hsize{\hfil 君子豹变。贵贱何常。青本奴虏。忽升戎行。\hfil}\linebreak
\hbox to \hsize{\hfil 姊配皇极。身尚平阳。宠荣斯僭。取乱彝章。\hfil}\linebreak
\hbox to \hsize{\hfil 嫖姚继踵。再静边方。\hfil}%
\par}
\end{document}
